% Standalone problem document, generated from the adjacent README.md.
% Compile with XeLaTeX. Mathematical content is maintained in Markdown.
\documentclass[11pt,a4paper]{article}
\usepackage[fontsize=10.5pt]{scrextend}
\usepackage[margin=22mm,headheight=15pt,headsep=9mm,footskip=11mm]{geometry}
\usepackage{fontspec}
\setmainfont{texgyrepagella}[Extension=.otf,UprightFont=*-regular,BoldFont=*-bold,ItalicFont=*-italic,BoldItalicFont=*-bolditalic]
\setsansfont{texgyreheros}[Extension=.otf,UprightFont=*-regular,BoldFont=*-bold,ItalicFont=*-italic,BoldItalicFont=*-bolditalic]
\setmonofont{texgyrecursor}[Extension=.otf,UprightFont=*-regular,BoldFont=*-bold,ItalicFont=*-italic,BoldItalicFont=*-bolditalic,Scale=MatchLowercase]
\usepackage{amsmath,amssymb}
\usepackage{unicode-math}
\setmathfont{texgyrepagella-math.otf}
\usepackage{xcolor,microtype,enumitem,needspace,fancyhdr}
\usepackage{bookmark}
\definecolor{ink}{HTML}{17344B}
\definecolor{accent}{HTML}{197D80}
\definecolor{muted}{HTML}{596773}
\hypersetup{colorlinks=true,linkcolor=accent,urlcolor=accent,pdftitle={IE-01 - Forsythe's
conjecture beyond restart length
two},pdfauthor={Open Problems in Numerical Linear Algebra}}
\urlstyle{same}
\setlength{\parindent}{0pt}
\setlength{\parskip}{5pt plus 1pt minus 1pt}
\linespread{1.04}
\setlength{\emergencystretch}{3em}
\widowpenalty=10000
\clubpenalty=10000
\displaywidowpenalty=10000
\allowdisplaybreaks[1]
\setlist{nosep,leftmargin=1.5em}
\providecommand{\tightlist}{\setlength{\itemsep}{0pt}\setlength{\parskip}{0pt}}
\providecommand{\pandocbounded}[1]{#1}
\pagestyle{fancy}
\fancyhf{}
\fancyhead[L]{\sffamily\footnotesize\color{muted}OPEN PROBLEMS IN NUMERICAL LINEAR ALGEBRA}
\fancyhead[R]{\sffamily\bfseries\footnotesize\color{accent}IE-01}
\fancyfoot[L]{\parbox[b]{0.9\textwidth}{\sffamily\fontsize{8}{10}\selectfont\color{muted}Editorial ratings; a bounded search cannot certify that a problem remains unsolved.\\Literature check: 2026-09-08}}
\fancyfoot[R]{\sffamily\footnotesize\color{muted}\thepage}
\renewcommand{\headrulewidth}{0.3pt}
\renewcommand{\headrule}{\hbox to\headwidth{\color{accent}\leaders\hrule height \headrulewidth\hfill}}
\makeatletter
\renewcommand{\section}{\Needspace{5\baselineskip}\@startsection{section}{1}{0pt}{12pt plus 2pt minus 2pt}{4pt}{\sffamily\bfseries\large\color{ink}}}
\renewcommand{\subsection}{\Needspace{4\baselineskip}\@startsection{subsection}{2}{0pt}{9pt plus 1pt minus 1pt}{3pt}{\sffamily\bfseries\normalsize\color{ink}}}
\makeatother
\setcounter{secnumdepth}{0}
\begin{document}
{\sffamily\small\color{muted}Linear systems and elimination\par}
\vspace{3pt}
{\raggedright\hyphenpenalty=10000\exhyphenpenalty=10000\sffamily\bfseries\fontsize{21}{24}\selectfont\color{ink}Forsythe's
conjecture beyond restart length two\par}
\vspace{7pt}
{\sffamily\small\textbf{Difficulty:} extreme\quad\textbf{Importance:} interesting
to the community\par}
\vspace{4pt}
{\color{accent}\hrule height 0.8pt}
\vspace{6pt}
\textbf{Status:} open; the case of restart length two has a recent proof
claim.

\subsection{Problem statement}\label{problem-statement}

Let \(A\in\mathbb R^{n\times n}\) be symmetric positive definite,
\(b,x_0\in\mathbb R^n\), and \(3\leq s<n\). At each restart, perform
exactly \(s\) exact-arithmetic conjugate-gradient steps, starting from
the last iterate, and discard the previous search directions. Denote the
iterate after restart cycle \(j\) by \(x_j\). Suppose this process never
terminates exactly, and put

\[
r_j=b-Ax_j,\qquad y_j=r_j/\|r_j\|_2.
\]

Prove or disprove that both \((y_{2j})_{j\geq0}\) and
\((y_{2j+1})_{j\geq0}\) converge in \(\mathbb R^n\), for every such
input and restart length. The question concerns directions, not whether
the residual norms tend to zero. Exact termination is excluded so every
normalization exists.

\subsection{References}\label{references}

Faber, Liesen, and Tichý,
\href{https://doi.org/10.1007/s10543-023-00991-x}{\emph{On the Forsythe
conjecture}}, BIT 63 (2023), §2, especially the displayed Forsythe
conjecture. Colbrook, Stepaniants, and Townsend,
\href{https://arxiv.org/abs/2608.02852}{\emph{A Proof of the Forsythe
Conjecture for the Two-Step Restarted Conjugate Gradient Method}},
August 2026, §1 and main theorem.

\subsection{Status check}\label{status-check}

Searches for \textit{Forsythe conjecture 2026 proof} found the August
preprint. Its claimed theorem covers \(s=2\), not all \(s\geq3\). The
surviving range is therefore stated explicitly; older claims that every
\(s\geq2\) remains open are obsolete.
\end{document}
